%#!latexmk -pdfdvi -f main.tex
%#LPR acroread %s

\RequirePackage{plautopatch}
\documentclass[paper=a4,
fontsize=12pt,
reqno,
uplatex,
dvipdfmx]{jlreq}
%
\jlreqsetup{
thebibliography_heading={\section*{\refname}},
}

%
%%%%%%%%%%%%%%% Include Package file %%%%%%%%%%%%%%%
\usepackage{amsmath,amssymb,amsthm}
%\usepackage{mathrsfs}
\usepackage{graphicx,color}
\usepackage{newtxtext}
\usepackage[varg]{newtxmath}
\usepackage[shortlabels]{enumitem}
\usepackage[haranoaji]{pxchfon}

\usepackage[alphabetic, y2k]{amsrefs}
%
%%%%%%%%%%%%%%% Paper size settings %%%%%%%%%%%%%%%
%
%
\usepackage{geometry}%
\geometry{left=20mm,right=20mm,top=15mm,bottom=25mm}
\geometry{a4paper}
%

%%%%%%%%%%%%%%% mathtools %%%%%%%%%%%%%%%
%
\usepackage{mathtools}
\mathtoolsset{showonlyrefs=true}
%
%
%
%%%%%%%%%%%%%%% hyperref %%%%%%%%%%%%%%%
%
\usepackage{hyperref}
\usepackage{pxjahyper}
%
%
%
%%%%%%%%%%%%%%% Setting of Theorem-like environment %%%%%%%%%%%%%%%
% [Italic font in Theorem-like environment]
%\theoremstyle{definition} % if change the roman font
\theoremstyle{definition} % Change the Roman font
%\newtheorem{definition}{定義}[chapter]
\newtheorem{theorem}{定理}
\newtheorem{proposition}[theorem]{命題}
\newtheorem{lemma}[theorem]{補題}
\newtheorem{corollary}[theorem]{系}
\newtheorem{problem}{問題}
\newtheorem{example}{例}
\newtheorem*{definition}{定義}
\newtheorem*{remark}{注意}
%
 
%
%%%%%%%%%% Proof %%%%%
\renewcommand{\proofname}{証明}  
% ((Use the environment in proof))
% \begin{proof}
% \end{proof}
% white box in the end of proof
%
%
% make new QED
% \QED{foo}
% we write a black box in the right and output foo after the black box
\newcommand{\QED}[1]{\hfill$\blacksquare$ #1\par}
% 
%
%
%%%%%%%%% 1.1 times Line stretch %%%%%%%%%
% If you do not use Japanese, that command should be commented out.
%
%
%%%%%%% Mathematical symbols %%%%%%%%%%%%%
%
\DeclareMathOperator{\supp}{supp}
\DeclareMathOperator{\diam}{diam}
\DeclareMathOperator{\sgn}{sgn}
\DeclareMathOperator{\Ker}{Ker}
\DeclareMathOperator{\Ran}{Ran}
\DeclareMathOperator{\dist}{dist}
\DeclareMathOperator{\tr}{tr}
\DeclareMathOperator{\Div}{div}
\DeclareMathOperator{\Int}{Int}
\DeclareMathOperator*{\osc}{osc}
\DeclareMathOperator{\Sym}{Sym}
\DeclareMathOperator*{\essosc}{ess.osc}
\DeclareMathOperator*{\esssup}{ess\,sup}

\newcommand{\N}{\mathbb{N}}
\newcommand{\Z}{\mathbb{Z}}
\newcommand{\Q}{\mathbb{Q}}
\newcommand{\R}{\mathbb{R}}
\newcommand{\Cplx}{\mathbb{C}}
\newcommand{\Sp}{\mathbb{S}}
\newcommand{\intmean}[1]{\,\rule[3pt]{10pt}{0.3pt}\hspace{-13pt}\int_{#1}}
% integral mean
% \intmean{D} is integral mean for D
\newcommand{\iintmean}[1]{\,\rule[3pt]{20pt}{0.3pt}\hspace{-22pt}\iint_{#1}}
% double integral mean
% \iintmean{D} is integral mean for D
%
%
%%%%%%%%%%%%%%% FRAME environment make a frame(similar to screen environment)
\newenvironment{FRAME}{\begin{trivlist}\item[]
 \hrule
 \hbox to \linewidth\bgroup
   \advance\linewidth by -30pt
   \hsize=\linewidth
   \vrule\hfill
   \vbox\bgroup
     \vskip15pt
     \def\thempfootnote{\arabic{mpfootnote}}% footnote is arabic number
     \begin{minipage}{\linewidth}}{%
     \end{minipage}\vskip15pt
   \egroup\hfill\vrule
 \egroup\hrule
\end{trivlist}}


%
%%%%%%%%%%%%%%% Counter setting in the equation's reference %%%%%%%%%%%%%%%
% For amsart, use this counter setting
%
%
%%% Definition of the \section (not use the \scshape)
%
%
%\usepackage{input/myamssec}
%
%%% \par after proofname
%
%\usepackage{input/myamspf}
%
%%% \par after theorem environment
%
%\usepackage{input/myamsthm}
%

%%%%%%%%%% Framed settings %%%%%
%
%\definecolor{shadecolor}{gray}{0.9}
%\newenvironment{shadedf}{%
%  \def\FrameCommand{\fcolorbox{white}{shadecolor}}%
%  \MakeFramed {\advance\hsize-\width \FrameRestore}%
%  }%
%{\endMakeFramed}
%
%
%


%%%%%%%%%%%%%%% temporary Package %%%%%%%%%%%%%%%

\usepackage[inline,right]{showlabels}   % Output of cite and reference

%\usepackage{showkeys}   % Output of cite and reference
%\usepackage{refcheck}   % check of no uselabel 

%%%%%%%%%%%%%%% title Setting %%%%%%%%%%%%%%%
%
\title{拡散方程式に対する自己相似変換}
\author{}
%
%
%%%%%%%%%%%%%%% page style setting %%%%%%%%%%%%%%%
%
\allowdisplaybreaks[1]
\DeclareMathOperator{\rot}{rot}
\let\Re\relax
\let\Im\relax
\DeclareMathOperator{\Re}{Re}
\DeclareMathOperator{\Im}{Im}

\usepackage{bm}
\renewcommand{\vec}[1]{\bm{#1}}


%
%%%%%%%%%%%%%%% end preamble %%%%%%%%%%%%%%%
%


\begin{document}
\maketitle

\begin{abstract}
 このノートは\cite{MR2656972}をもとにしたものである．
\end{abstract}


$n$ 次元 Euclid 空間 $\mathbb{R}^n$ における拡散方程式
\begin{equation}
 \label{eq:DiffusionEquation}
 \frac{\partial u}{\partial t}(x, t)
  -
  \Delta u(x, t) = 0, 
  \quad x \in \mathbb{R}^n,
  \, t > 0 \tag{D}
\end{equation}
を考える．$\lambda > 0$ に対し、スケール変換
\begin{equation}
 \label{eq:Self-Similar-scale-transform}
 \tilde{u}_{\lambda}(\tilde{x}, \tilde{t})
  =
  \lambda^{-\frac{n}{2}} u(x, t), 
  \quad 
  \tilde{x} = \lambda x, \, \tilde{t} = \lambda^2 t
\end{equation}
を導入する．

\begin{proposition}
 $u: \mathbb{R}^n \times (0, \infty) \rightarrow \mathbb{R}$と$\lambda >
 0$ に対して、$\tilde{u}_{\lambda}$
 を~\eqref{eq:Self-Similar-scale-transform}で定める．
 \begin{enumerate}
  \item 任意の $\lambda > 0$ に対して、次が成り立つ: 
   \begin{equation}
    \int_{\mathbb{R}^n} 
     u(x, t)
     \, dx 
     = 
     \int_{\mathbb{R}^n}
     \tilde{u}_{\lambda}(\tilde{x}, \tilde{t})
     \, 
     d\tilde{x}
   \end{equation}
  \item $u$ が \eqref{eq:DiffusionEquation} を満たすならば、任意の
	$\lambda > 0$ に対して，$\tilde{u}_{\lambda}$ も次の拡散方程式の
	解となる: 
	\begin{equation}
	 \frac{\partial \tilde{u}_{\lambda}}{\partial \tilde{t}}
	  (\tilde{x}, \tilde{t}) 
	  -
	  \Delta_{\tilde{x}} \tilde{u}_{\lambda}
	  (\tilde{x}, \tilde{t})
	  =
	  0,
	  \quad
	  \tilde{x} \in \mathbb{R}^n, \, 
	  \tilde{t} > 0.
	\end{equation}
 \end{enumerate}
\end{proposition}

さて，$u$ を \eqref{eq:DiffusionEquation} の解とするとき，全ての
$\lambda > 0$ に対して
\begin{equation}
 \lambda^{-\frac{n}{2}}
  u
  \left(
   \frac{x}{\lambda}, 
   \frac{t}{\lambda^2}
  \right)
  =
  u(x, t), \quad \forall x \in \mathbb{R}^n, \, t > 0
\end{equation}
となる関数があったとすると(自己相似解という)，$\lambda =
\frac{1}{\sqrt{t}}$ とおくことで
\begin{equation}
 t^{-\frac{n}{2}} u\left(\frac{x}{\sqrt{t}}, 1\right) = u(x, t), \quad x \in \mathbb{R}^n, \, t > 0
\end{equation}
がわかる．そこで，新しい変数 $y = \frac{x}{\sqrt{t}}$ を導入して
\begin{equation}
 u(x, t)
  =
  t^{-\frac{n}{2}}
  \tilde{v}(y, t),
  \quad
  y = \frac{x}{\sqrt{t}}
\end{equation}
とおく．

\begin{proposition}
 $u: \mathbb{R}^n \times (0, \infty) \rightarrow \mathbb{R}$
 に対し，
 $\tilde{v}$ を上記で定める．このとき，全ての $t > 0$ に対して次が成り立つ:
 \begin{equation}
  \int_{\mathbb{R}^n} u(x, t) \, dx 
   =
   \int_{\mathbb{R}^n} \tilde{v}(y, t) \, dy.
 \end{equation}
\end{proposition}

次に，$u$ が \eqref{eq:DiffusionEquation} を満たすときの $\tilde{v}$ の
満たす方程式を導出しよう．
\begin{align}
 \frac{\partial u}{\partial t}(x, t)
 &=
 \frac{\partial}{\partial t}
 \left(
 t^{-\frac{n}{2}} \tilde{v}\left(\frac{x}{\sqrt{t}}, t\right) \right) \tag{592} \\
 &= -\frac{n}{2} t^{-\frac{n}{2}-1} \tilde{v}(y, t) - \frac{1}{2} t^{-\frac{n}{2}-1} y \cdot (\nabla_y \tilde{v})(y, t) + t^{-\frac{n}{2}} \frac{\partial \tilde{v}}{\partial t}(y, t),
\end{align}
かつ
\begin{equation}
 \Delta_x u(x, t)
  = 
  \Delta_x
  \left(
   t^{-\frac{n}{2}}
   \tilde{v}\left(\frac{x}{\sqrt{t}}, t\right) 
  \right)
  =
  t^{-\frac{n}{2}-1} 
  (\Delta_y \tilde{v})(y, t)
\end{equation}
を~\eqref{eq:DiffusionEquation}に代入して整理すると
\begin{equation}
 t \frac{\partial \tilde{v}}{\partial t}(y, t)
  -
  (\Delta_y \tilde{v})(y, t)
  -
  \frac{1}{2} y \cdot \nabla_y \tilde{v}(y, t)
  -
  \frac{n}{2} \tilde{v}(y, t) = 0 
\end{equation}
となる．ここで，
\begin{equation}
 \Div_y
  \left(
   \frac{1}{2}
   y
   \tilde{v}(y, t)
  \right)
  =
  \frac{n}{2}
  \tilde{v}(y, t)
  + 
  \frac{1}{2}
  y
  \cdot
  \nabla_y
  \tilde{v}(y, t)
\end{equation}
に注目すると，
\begin{equation}
 t \frac{\partial \tilde{v}}{\partial t}(y, t)
  -
  (\Delta_y \tilde{v})(y, t)
  -
  \Div_y 
  \left(
   \frac{1}{2}
   \tilde{v}(y, t)
   y
  \right)
  = 0
\end{equation}
が得られる．

さらに，時間変数を変換するために $s = \log t$ という変数変換を導入する．
\begin{equation}
 v(y, s)
  =
  \tilde{v}(y, t), 
  \quad s = \log t
\end{equation}
とおくと，
\begin{equation}
 \frac{\partial \tilde{v}}{\partial t}(y, t)
  =
  \frac{\partial}{\partial t} v(y, \log t)
  = 
  \frac{\partial v}{\partial s}(y, s) \frac{1}{t}
\end{equation}
となるから
\begin{equation}
 \frac{\partial v}{\partial s}(y, s)
  -
   (\Delta_y v)(y, s)
   - 
   \text{div}_y 
   \left(
    \frac{1}{2} v(y, s) y
   \right)
   =
   0, \quad y \in \mathbb{R}^n, \, s \in \mathbb{R}
\end{equation}
が得られる．

\begin{definition}
 $u: \mathbb{R}^n \times (0, \infty) \rightarrow \mathbb{R}$ に対し、
 \begin{equation}
  u(x, t)
   = 
   t^{-\frac{n}{2}}
   v(y, s), 
   \quad 
   y = \frac{x}{\sqrt{t}}, \, s = \log t
 \end{equation}
 で定まる変換を，~\eqref{eq:DiffusionEquation} に関する\textbf{前向自己
 相似変換（forward self-similar transformation）}という．
\end{definition}

\begin{proposition}
 $u: \mathbb{R}^n \times (0, \infty) \rightarrow \mathbb{R}$ に対し，
 次が成り立つ．
 \begin{enumerate}
  \item 全ての $t > 0$ に対し、
	\begin{equation}
	 \int_{\mathbb{R}^n} u(x, t) \, dx = \int_{\mathbb{R}^n} v(y, s) \, dy
	\end{equation}
  \item $u$ が \eqref{eq:DiffusionEquation} を満たすならば，
	\begin{equation}
	 \frac{\partial v}{\partial s}(y, s)
	  -
	  (\Delta_y v)(y, s)
	  -
	  \Div_y
	  \left(
	   \frac{1}{2} v(y, s) y
	  \right)
	  =
	  0, 
	  \quad y \in \mathbb{R}^n, 
	  \, s \in \mathbb{R}
	\end{equation}
	が成り立つ．
 \end{enumerate}
\end{proposition}

さて，$u$ を正値とすると，前向自己相似変換 $v$ もまた正値となる．
このとき，方程式は次のように変形できる: 
\begin{align}
 \frac{\partial v}{\partial s}(y,s)
 &=
 \Delta v(y,s)
 +
 \text{div} 
 \left(
 \frac{1}{2} v(y, s) y
 \right) \\
 &= \Div
 \left(
 \nabla v(y, s)
 + 
 \frac{1}{2} v(y, s) y
 \right) \\
 &= \Div
 \left(
 v(y, s)
 \nabla
 \left( 
 \log v(y, s) + \frac{1}{4} |y|^2
 \right) 
 \right)
\end{align}
天下りではあるが，次のエネルギーの微分を計算してみる:
\begin{align}
 &\quad
 \frac{d}{ds}
  \int_{\mathbb{R}^n}
  \left( 
  (\log v(y, s) - 1)
  + 
  \frac{1}{4} |y|^2
  \right)
  v(y, s)
  \, dy
 \\
  &=
  \int_{\mathbb{R}^n} 
  \left(
  \log v(y, s) + \frac{1}{4} |y|^2 
  \right)
  \frac{\partial v}{\partial s}(y,s)
  \, dy
  \\
  &=
  \int_{\mathbb{R}^n}
  \left(
  \log v(y, s) + \frac{1}{4} |y|^2 
  \right) 
  \Div
  \left(
  v(y,s)
  \nabla 
  \left(
  \log v(y,s)
  + 
  \frac{1}{4} |y|^2
  \right)
  \right)
  \,dy
 \end{align}
これに部分積分の公式を用いると，$v$ が無限遠で十分に早く減衰していれば，
境界項は $0$ となるため，
\begin{equation}
 \frac{d}{ds} 
  \int_{\mathbb{R}^n}
  \left(
   (\log v(y, s) - 1) + \frac{1}{4} |y|^2
  \right)
  v(y, s)
  \, dy 
  =
  -\int_{\mathbb{R}^n}
  \left|
   \nabla 
   \left( 
    \log v(y, s) + \frac{1}{4} |y|^2
   \right) 
  \right|^2 v(y, s)
  \, dy
\end{equation}
となる．さらに，時間変数について積分することで，次の自然境界条件の下での
エネルギー等式が得られる:
\begin{multline}
 \int_{\mathbb{R}^n}
 \left( 
 (\log v(y, s) - 1) + \frac{1}{4} |y|^2
 \right) v(y, s)
 \, dy
 +
 \int_{0}^s
 \left(
 \int_{\mathbb{R}^n}
 \left|
 \nabla 
 \left(
 \log v(y, \tau)
 + 
 \frac{1}{4} |y|^2 
 \right)
 \right|^2
 v(y, \tau)
 \, dy 
 \right)
 d\tau \\
 = 
 \int_{\mathbb{R}^n}
 \left( 
 (\log v(y, 0) - 1) + \frac{1}{4} |y|^2
 \right)
 v(y, 0)
 \, dy 
\end{multline}
このことから，$s \rightarrow \infty$ としたとき，
\begin{equation}
 \nabla
  \left(
   \log v(y, s)
   +
   \frac{1}{4} |y|^2
  \right) 
  \rightarrow 
  0
\end{equation}
となることが示唆される．したがって，$\log v(y, s) + \frac{1}{4} |y|^2$
が定数に収束すればよいことから，
\begin{equation}
 v(y, s) 
  \rightarrow
  C e^{-\frac{1}{4} |y|^2}
  \quad 
  s \rightarrow \infty
\end{equation}
となることが推察される．

\begin{theorem}
 $u$ を \eqref{eq:DiffusionEquation} の解とし，$v$ を $u$ の前向自己相似
 変換とする．定数 $C > 0$ に対し，
 \begin{equation}
  \| 
   v(y, s)
   -
   C e^{-\frac{|y|^2}{4}}
   \|_{L^p(\mathbb{R}_y^n)}
   =
   \phi(s)
 \end{equation}
 とおく．このとき，$t>0$に対して
 \begin{equation}
  \| 
   u(x, t)
   -
   C t^{-\frac{n}{2}} e^{-\frac{|x|^2}{4t}}
   \|_{L^p(\mathbb{R}_x^n)}
   =
   t^{-\frac{n}{2}
   \left(1 - \frac{1}{p}
   \right)}
   \phi(\log t)
 \end{equation}
 が成り立つ．
\end{theorem}

\begin{proof}
 自己相似変換の定義より，
 \begin{align}
  \| 
  v(y, s)
  -
  C e^{-\frac{|y|^2}{4}}
  \|_{L^p(\mathbb{R}_y^n)}^p
  &=
  \int_{\mathbb{R}^n}
  \left|
  v(y, s)
  -
  C e^{-\frac{|y|^2}{4}}
  \right|^p
  \, dy \\
  &= 
  \int_{\mathbb{R}^n}
  \left| 
  t^{\frac{n}{2}} u(x, t) - C e^{-\frac{|y|^2}{4}}
  \right|^p
  \, dy \\
  &=
  \int_{\mathbb{R}^n} t^{\frac{np}{2}} 
  \left| 
  u(x, t)
  -
  C t^{-\frac{n}{2}} e^{-\frac{|y|^2}{4}}
  \right|^p
  \, dy
 \end{align}
 となる，ここで $y = \frac{x}{\sqrt{t}}$，$s = \log t$ と変数変換すると，
 $dy = t^{-\frac{n}{2}} dx$ であるから，
 \begin{align}
  \| 
  v(y, s)
  -
  C e^{-\frac{|y|^2}{4}} 
  \|_{L^p(\mathbb{R}_y^n)}^p 
  &=
  t^{\frac{np}{2}} 
  t^{-\frac{n}{2}}
  \int_{\mathbb{R}^n}
  \left|
  u(x, t) - C t^{-\frac{n}{2}} e^{-\frac{|x|^2}{4t}}
  \right|^p
  \, dx \\
  &=
  t^{\frac{n}{2}(p - 1)} 
  \int_{\mathbb{R}^n}
  \left| 
  u(x, t) - C t^{-\frac{n}{2}} e^{-\frac{|x|^2}{4t}}
  \right|^p
  \, dx
 \end{align}
 となる．両辺を $\frac{1}{p}$ 乗して整理すれば，
 \begin{equation}
  \|
   u(x, t) 
   -
   C t^{-\frac{n}{2}} e^{-\frac{|x|^2}{4t}}
   \|_{L^p(\mathbb{R}_x^n)}
   = 
   t^{-\frac{n}{2}\left(1 - \frac{1}{p}\right)} 
   \|
   v(y, s) - C e^{-\frac{|y|^2}{4}}
   \|_{L^p(\mathbb{R}_y^n)}
   = 
   t^{-\frac{n}{2}\left(1 - \frac{1}{p}\right)} \phi(\log t) 
 \end{equation}
 が得られる．
\end{proof}

\section{応用}
もし，ある定数 $C_1 > 0$ と $\lambda > 0$ に対して
\begin{equation}
 \|
  v(y, s) - C e^{-\frac{|y|^2}{4}}
  \|_{L^p(\mathbb{R}_y^n)}
  \le
  C_1 e^{-\lambda s} 
\end{equation}
なる評価が得られたとすると，
\begin{equation}
 \|
  u(x, t) - C t^{-\frac{n}{2}} e^{-\frac{|x|^2}{4t}}
  \|_{L^p(\mathbb{R}_x^n)}
  \le
  C_1 t^{-\frac{n}{2}
  \left(
   1 - \frac{1}{p}
  \right)}
  e^{-\lambda \log t}
  =
  C_1 t^{-\frac{n}{2}
  \left(
   1 - \frac{1}{p}
  \right)
  - \lambda}
\end{equation}
となる．つまり，$u(x, t)$ が十分大きな時間において Gauss 核（基本解）に
漸近すること，およびその剰余項が代数的なオーダー $t^{-\lambda}$ で減衰す
ることがわかる．



% bibtex Settings
\bibliography{references}
%\bibliographystyle{myjalpha} % in Japanese

%\bibliographystyle{amsplain} % not in Japanese


\end{document}




\end{document}